\section{推荐阅读}

\begin{readingbox}
\textbf{基础}
\begin{itemize}[nosep]
  \item John Holland, \emph{Emergence}.
  \item \emph{Introduction to the Theory of Complex Systems}.
  \item \emph{Complexity Science}，用于扩展关于自组织与多尺度秩序的系统词汇。
\end{itemize}
\textbf{生物学与综述层}
\begin{itemize}[nosep]
  \item \emph{Metastability demystified: the foundational past, the pragmatic
        present and the promising future} （Nature Reviews Neuroscience, 2024,
        medium）。
  \item \emph{Associative conditioning in gene regulatory network models
        increases integrative causal emergence} （Communications Biology, 2025,
        medium）。
  \item \emph{Cortical connectivity, local dynamics and stability correlates of
        global conscious states} （Communications Biology, 2025, 作为生物学背景属 strong，作为本章迁移主张属 medium）。
\end{itemize}
\textbf{数学桥梁}
\begin{itemize}[nosep]
  \item Eugene Izhikevich, \emph{Dynamical Systems in Neuroscience}.
  \item 关于大脑中吸引子与积分器网络的综述文献
        （Nature Reviews Neuroscience, 2022, medium）。
\end{itemize}
\textbf{阅读纪律}
\begin{itemize}[nosep]
  \item 用生物学与综述层来支撑“有组织状态形成与稳定性”的正当性。
  \item 用系统与数学层来澄清：哪一种高层模式才算真正具有涌现性质。
  \item 不要把前沿意识物理学方案当作本章核心主张的证据。
\end{itemize}
\end{readingbox}